\subsection{单 g-模的维数}

若 \(\mu \in \mathfrak{h}^*\)，置 \(\mathrm{J}(e^{\mu}) = \sum_{w\in \mathrm{W}}\varepsilon (w)e^{w\mu}\)，参见~第 VI 章第3节第3号。

定理 2。令 E 为有限维单 g-模，\(\lambda\) 为其最高权，\((\cdot|\cdot)\) 为 \(h^{*}\) 上 W-不变的非退化正对称双线性型。则：

\[
\dim \mathrm{E} = \prod_ {\alpha \in \mathrm{R} _ {+}} \frac {\langle \lambda + \rho , H _ {\alpha} \rangle}{\langle \rho , H _ {\alpha} \rangle} = \prod_ {\alpha \in \mathrm{R} _ {+}} \left(1 + \frac {(\lambda | \alpha)}{(\rho | \alpha)}\right).
\]

令 \(\mathrm{T}\) 为不定元。对所有 \(\nu \in \mathrm{P}\)，以 \(f_{\nu}\) 表示从 \(\mathbf{Z}[\mathrm{P}]\) 到 \(\mathbf{R}[[\mathrm{T}]]\) 的同态，它对所有 \(e^{\mu}\) 将 \(e^{(\nu |\mu)\mathrm{T}}\) 映为 \(\mu \in \mathrm{P}\)。则 \(\dim \mathrm{E}\) 是级数 \(f_{\nu}(\operatorname{ch}\mathrm{E})\) 的常数项。

对所有 \(\mu, \nu \in \mathrm{P}\)，有

\[
\begin{array}{c} f _ {\nu} (\mathrm{J} (e ^ {\mu})) = \sum_ {w \in \mathrm{W}} \varepsilon (w) e ^ {(\nu | w \mu) \mathrm{T}} \\ = \sum_ {w \in \mathrm{W}} \varepsilon (w) e ^ {(w ^ {- 1} \nu | \mu) \mathrm{T}} = f _ {\mu} (\mathrm{J} (e ^ {\nu})). \end{array}
\]

特别地，根据第 VI 章第3节第3号的公式 (3)，

\[
f _ {\rho} (\mathrm{J} (e ^ {\mu})) = f _ {\mu} (\mathrm{J} (e ^ {\rho})) = e ^ {(\mu | \rho) \mathrm{T}} \prod_ {\alpha \in \mathrm{R} _ {+}} (1 - e ^ {- (\mu | \alpha) \mathrm{T}}).
\]

因此，置 \(\mathrm{Card}(\mathrm{R}_{+})=\mathrm{N}\)，

\[
f _ {\rho} (\mathrm{J} (e ^ {\mu})) \equiv \mathrm{T} ^ {\mathrm{N}} \prod_ {\alpha \in \mathrm{R}} (\mu | \alpha) \pmod {\mathrm{T} ^ {\mathrm{N} + 1} \mathbf {R} [ [ \mathrm{T} ] ]}.
\]

等式 \(\mathrm{J}(e^{\lambda +\rho}) = \mathrm{ch}(\mathrm{E}).\mathrm{J}(e^{\rho})\)（定理 1）从而蕴含

\[
\mathrm{T} ^ {\mathrm{N}} \prod_ {\alpha \in \mathrm{R} _ {+}} (\lambda + \rho | \alpha) \equiv f _ {\rho} (\operatorname{ch} \mathrm{E}). \mathrm{T} ^ {\mathrm{N}} \prod_ {\alpha \in \mathrm{R} _ {+}} (\rho | \alpha) \pmod {\mathrm{T} ^ {\mathrm{N} + 1} \mathbf {R} [ [ \mathrm{T} ] ]}
\]

因此

\[
\dim \mathrm{E} = \left(\prod_ {\alpha \in \mathrm{R} _ {+}} (\lambda + \rho | \alpha)\right) / \left(\prod_ {\alpha \in \mathrm{R} _ {+}} (\rho | \alpha)\right) = \prod_ {\alpha \in \mathrm{R} _ {+}} \left(1 + \frac {(\lambda | \alpha)}{(\rho | \alpha)}\right).
\]

现在，若 \(\alpha \in \mathrm{R}_{+}\)，则可将 \(\alpha\) 识别为 \(\mathfrak{h}_{\mathbf{R}}\) 中与 \(H_{\alpha}\) 成比例的元素，因而

\[
(\lambda + \rho | \alpha) / (\rho | \alpha) = \langle \lambda + \rho , H _ {\alpha} \rangle / \langle \rho , H _ {\alpha} \rangle .
\]

例：1）在第1号的例子中，我们得到

\[
\dim \mathrm{V} (m) = \left(\frac {m}{2} \alpha + \frac {\alpha}{2}\right) (H _ {\alpha}) / \frac {\alpha}{2} (H _ {\alpha}) = m + 1,
\]

这正是我们在 §1 中已经知道的结果。

\begin{enumerate}
\def\labelenumi{\arabic{enumi})}
\setcounter{enumi}{1}
\tightlist
\item
  取 g 为 \(G_{2}\) 型的可分裂单李代数，并采用第 VI 章图版 IX 的记号。在 \(h_{R}^{*}\) 上取 W-不变的正对称型 \((\cdot|\cdot)\)，使得 \((\alpha_{1}|\alpha_{1})=1\)。则 \(\rho=\varpi_{1}+\varpi_{2}\)，且
\end{enumerate}

\[
\begin{array}{c} (\varpi_ {1} | \alpha_ {1}) = \frac {1}{2}, (\varpi_ {1} | \alpha_ {2}) = 0, (\varpi_ {1} | \alpha_ {2} + \alpha_ {1}) = \frac {1}{2}, \\ (\varpi_ {1} | \alpha_ {2} + 2 \alpha_ {1}) = 1, (\varpi_ {1} | \alpha_ {2} + 3 \alpha_ {1}) = \frac {3}{2}, (\varpi_ {1} | 2 \alpha_ {2} + 3 \alpha_ {1}) = \frac {3}{2}, \\ (\varpi_ {2} | \alpha_ {1}) = 0, (\varpi_ {2} | \alpha_ {2}) = \frac {3}{2}, (\varpi_ {2} | \alpha_ {2} + \alpha_ {1}) = \frac {3}{2}, \\ (\varpi_ {2} | \alpha_ {2} + 2 \alpha_ {1}) = \frac {3}{2}, (\varpi_ {2} | \alpha_ {2} + 3 \alpha_ {1}) = \frac {3}{2}, (\varpi_ {2} | 2 \alpha_ {2} + 3 \alpha_ {1}) = 3. \end{array}
\]

因此，若 \(n_1, n_2\) 是 \(\geq 0\) 的整数，则最高权为 \(n_1\varpi_1 + n_2\varpi_2\) 的单表示的维数为

\[
\begin{array}{c} \left(1 + \frac {n _ {1} / 2}{\frac {1}{2}}\right) \left(1 + \frac {3 n _ {2} / 2}{\frac {3}{2}}\right) \left(1 + \frac {n _ {1} / 2 + 3 n _ {2} / 2}{\frac {1}{2} + \frac {3}{2}}\right) \left(1 + \frac {n _ {1} + 3 n _ {2} / 2}{1 + \frac {3}{2}}\right) \\ \times \left(1 + \frac {3 n _ {1} / 2 + 3 n _ {2} / 2}{\frac {3}{2} + \frac {3}{2}}\right) \left(1 + \frac {3 n _ {1} / 2 + 3 n _ {2}}{\frac {3}{2} + 3}\right) \\ = (1 + n _ {1}) (1 + n _ {2}) \left(1 + \frac {n _ {1} + 3 n _ {2}}{4}\right) \left(1 + \frac {2 n _ {1} + 3 n _ {2}}{5}\right) \left(1 + \frac {n _ {1} + n _ {2}}{2}\right) \\ \times \left(1 + \frac {n _ {1} + 2 n _ {2}}{3}\right) \\ = \frac {(1 + n _ {1}) (1 + n _ {2}) (2 + n _ {1} + n _ {2}) (3 + n _ {1} + 2 n _ {2}) (4 + n _ {1} + 3 n _ {2}) (5 + 2 n _ {1} + 3 n _ {2})}{5 !}. \end{array}
\]

特别地，最高权为 \(\varpi_{1}\)（分别为 \(\varpi_{2}\)）的基本表示的维数为 7（分别为 14）。

